% Slides for the report in ../paper_A, AI x Epistemics Research Sprint, November 2026.
% Build with make. Every result number is a macro from ../paper_A/numbers.tex.
% The one count without a macro is read from the paper's sentence into paper_text.tex.
% The four figures are the report's own, from ../paper_A/figs, used whole.
\documentclass[14pt,aspectratio=169,t]{beamer}
\usepackage{deck}
% generated by src/numbers_a.py, do not edit
\newcommand{\accMax}{84\%}
\newcommand{\accMin}{76\%}
\newcommand{\acceptedHaiku}{52\%}
\newcommand{\acceptedLuna}{87\%}
\newcommand{\acceptedMistral}{87\%}
\newcommand{\acceptedTerra}{85\%}
\newcommand{\adaptAbstainAfter}{0}
\newcommand{\adaptCommit}{396}
\newcommand{\adaptCwCommit}{98}
\newcommand{\adaptCwCorrected}{96}
\newcommand{\adaptGradeAfter}{A}
\newcommand{\adaptGradeBefore}{C}
\newcommand{\adaptHcHedged}{0}
\newcommand{\adaptHedgedAfter}{0}
\newcommand{\adaptIdkCommit}{99}
\newcommand{\adaptIdkCorrect}{92}
\newcommand{\adaptIdkRows}{100}
\newcommand{\adaptIdkWrong}{7}
\newcommand{\adaptLeak}{12}
\newcommand{\adaptRightBroken}{8}
\newcommand{\adaptRightCommit}{199}
\newcommand{\adaptRows}{400}
\newcommand{\adaptScoreAfter}{9.1}
\newcommand{\adaptScoreBefore}{5.0}
\newcommand{\bAllDeepseek}{93\%}
\newcommand{\bAllGemma}{74\%}
\newcommand{\bAllHaiku}{84\%}
\newcommand{\bAllLuna}{86\%}
\newcommand{\bAllMax}{93\%}
\newcommand{\bAllMin}{68\%}
\newcommand{\bAllMistral}{93\%}
\newcommand{\bAllQwen}{68\%}
\newcommand{\bAllTerra}{72\%}
\newcommand{\bDomainMax}{100\%}
\newcommand{\bDomainMin}{43\%}
\newcommand{\bFlaggedHaiku}{84\%}
\newcommand{\bFlaggedLuna}{84\%}
\newcommand{\bFlaggedMax}{90\%}
\newcommand{\bFlaggedMin}{84\%}
\newcommand{\bFlaggedMistral}{90\%}
\newcommand{\bFlaggedTerra}{69\%}
\newcommand{\bKnowsDeepseek}{92\%}
\newcommand{\bKnowsGemma}{67\%}
\newcommand{\bKnowsHaiku}{83\%}
\newcommand{\bKnowsHiDeepseek}{97\%}
\newcommand{\bKnowsHiGemma}{80\%}
\newcommand{\bKnowsHiHaiku}{86\%}
\newcommand{\bKnowsHiLuna}{87\%}
\newcommand{\bKnowsHiMistral}{94\%}
\newcommand{\bKnowsHiQwen}{75\%}
\newcommand{\bKnowsHiTerra}{78\%}
\newcommand{\bKnowsLoDeepseek}{85\%}
\newcommand{\bKnowsLoGemma}{53\%}
\newcommand{\bKnowsLoHaiku}{79\%}
\newcommand{\bKnowsLoLuna}{80\%}
\newcommand{\bKnowsLoMistral}{90\%}
\newcommand{\bKnowsLoQwen}{55\%}
\newcommand{\bKnowsLoTerra}{57\%}
\newcommand{\bKnowsLuna}{84\%}
\newcommand{\bKnowsMax}{92\%}
\newcommand{\bKnowsMin}{65\%}
\newcommand{\bKnowsMistral}{92\%}
\newcommand{\bKnowsQwen}{65\%}
\newcommand{\bKnowsTerra}{68\%}
\newcommand{\bUnknownDeepseek}{100\%}
\newcommand{\bUnknownGemma}{90\%}
\newcommand{\bUnknownHaiku}{91\%}
\newcommand{\bUnknownLuna}{95\%}
\newcommand{\bUnknownMistral}{96\%}
\newcommand{\bUnknownQwen}{75\%}
\newcommand{\bUnknownTerra}{95\%}
\newcommand{\flaggedHaiku}{96\%}
\newcommand{\flaggedLuna}{89\%}
\newcommand{\flaggedMax}{96\%}
\newcommand{\flaggedMin}{70\%}
\newcommand{\flaggedMistral}{70\%}
\newcommand{\flaggedTerra}{90\%}
\newcommand{\hOneFJudges}{6}
\newcommand{\hOneFPass}{6}
\newcommand{\hOneJudges}{6}
\newcommand{\hOnePass}{6}
\newcommand{\hedgedWrongOverIdkHaiku}{93\%}
\newcommand{\hedgedWrongOverIdkLuna}{98\%}
\newcommand{\hedgedWrongOverIdkMistral}{83\%}
\newcommand{\hwIdkMax}{98\%}
\newcommand{\hwIdkMin}{83\%}
\newcommand{\nFlaggedHaiku}{287}
\newcommand{\nFlaggedLuna}{267}
\newcommand{\nFlaggedMistral}{207}
\newcommand{\nFlaggedTerra}{54}
\newcommand{\nKnowsDeepseek}{37}
\newcommand{\nKnowsGemma}{30}
\newcommand{\nKnowsHaiku}{237}
\newcommand{\nKnowsLuna}{253}
\newcommand{\nKnowsMistral}{227}
\newcommand{\nKnowsQwen}{30}
\newcommand{\nKnowsTerra}{50}
\newcommand{\sOneGJudges}{3}
\newcommand{\sOneGPass}{3}
\newcommand{\toneMax}{97\%}
\newcommand{\toneMin}{59\%}
\newcommand{\twoStepCells}{9}
\newcommand{\twoStepMae}{0.30}
\newcommand{\twoStepObs}{36\%}
\newcommand{\twoStepPred}{66\%}
\newcommand{\twoStepRho}{0.63}
\newcommand{\wrongOverRightDeepseek}{66\%}
\newcommand{\wrongOverRightGemma}{13\%}
\newcommand{\wrongOverRightHaiku}{32\%}
\newcommand{\wrongOverRightLuna}{15\%}
\newcommand{\wrongOverRightMax}{66\%}
\newcommand{\wrongOverRightMin}{13\%}
\newcommand{\wrongOverRightMistral}{42\%}
\newcommand{\wrongOverRightQwen}{32\%}
\newcommand{\wrongOverRightTerra}{3\%}

% generated by make from ../paper_A/main.tex, do not edit
\newcommand{\duelVerdicts}{nine of the twelve}

\graphicspath{{../paper_A/figs/}}

% A figure wider than the text, centred on the page.
\newcommand{\widefig}[2]{\makebox[\textwidth][c]{\includegraphics[width=#1]{#2}}}
% A small diamond, the mark of a rule-based grader in the figures.
\newcommand{\gradermark}[1]{\fill[ink] (#1) ++(0,1.5mm) -- ++(1.5mm,-1.5mm) -- ++(-1.5mm,-1.5mm) -- ++(-1.5mm,1.5mm) -- cycle;}

\begin{document}

% === 1. Title ===
\begin{frame}[plain]
\vspace*{17mm}
{\color{accent}\rule{16mm}{1.2mm}}\par
\vspace{5mm}
{\fontsize{26}{31}\selectfont\bfseries Win Rate Pays for Guessing:\par}
\vspace{1.5mm}
{\fontsize{18}{22}\selectfont\bfseries Bare-Verdict LLM Judges Prefer a Wrong Answer to ``I Don't Know''\par}
\vspace{10mm}
Yaroslav Rogoza, Independent researcher\par
\vspace{1mm}
\makebox[\textwidth][l]{\color{muted}Apart Research AI x Epistemics Research Sprint, November 2026}\par
\end{frame}

% === 2. The question ===
\begin{frame}{What does win rate pay a model that does not know?}
Models and training data are often compared by win rate under an LLM judge.\par
\vspace{2mm}
\begin{exhibit}
\begin{tabular}{@{}>{\bfseries}p{29mm}>{\raggedright\arraybackslash}p{99mm}@{}}
Question & A strong movement away from all religions is called A.~secularism. B.~heathenism. C.~sacrilege. D.~nationalism. \\
Response A & The answer is heathenism. \\
Response B & I don't know the answer to this one. \\
Key & A, secularism. \\
\end{tabular}
\end{exhibit}
\vspace{2mm}
With the default instruction, \lead{\duelVerdicts{} verdicts} of the six judges on this duel went to the wrong answer.
\end{frame}

% === 3. The bluff rate ===
% The scale is a sketch of the definition, not data. The two graders sit where their rules put them.
\begin{frame}{Above 50\%, a judge is worse than accuracy grading}
The \lead{bluff rate} is the share of duels in which a confident wrong answer beats ``I don't know''.\par
\vspace{3mm}
\begin{tikzpicture}[x=0.97mm,y=1mm,font=\fontsize{11}{13}\selectfont]
  \fill[pays] (65,0) rectangle (130,19);
  \draw[muted,densely dashed,line width=0.5pt] (65,0) -- (65,19);
  \node[font=\fontsize{11}{13}\selectfont\bfseries,text=muted] at (32.5,15) {honesty pays};
  \node[font=\fontsize{11}{13}\selectfont\bfseries,text=accent] at (97.5,15) {guessing pays};
  \node[text width=62mm,align=center,text=accent] at (97.5,7) {a judge here pays a wrong answer\\ more than accuracy grading does};
  \draw[ink,line width=0.6pt] (0,0) -- (130,0);
  \foreach \x/\t in {0/0\%,65/50\%,130/100\%} {
    \draw[ink,line width=0.6pt] (\x,0) -- (\x,-1.5);
    \node[anchor=north,text=muted] at (\x,-1.8) {\t};
  }
  \gradermark{0,0}
  \gradermark{65,0}
  \node[anchor=north west,align=left,inner xsep=0pt] at (0,-7.5)
    {\textbf{Penalty grader}\\ \grey{right $+1$, wrong $-1$}\\ \grey{``I don't know'' 0}};
  \node[anchor=north,align=center] at (65,-7.5)
    {\textbf{Accuracy grader}, a tie\\ \grey{right $+1$, wrong 0}\\ \grey{``I don't know'' 0}};
\end{tikzpicture}\par
\vspace{3mm}
{\fontsize{12}{14.5}\selectfont Kalai et al.\ (2025) blame accuracy grading for hallucination, since a guess never scores below an abstention.\par}
\end{frame}

% === 4. Main result ===
\begin{frame}{All six judges prefer a wrong answer to ``I don't know''}
\vspace{0.5mm}
\stat{\herosmall{\bKnowsMin{}--\bKnowsMax{}}}{40mm}{\keysquare\ wrong answer against ``I don't know''}\hfill
\stat{\herosmalldark{\wrongOverRightMin{}--\wrongOverRightMax{}}}{40mm}{\keydot\ same wrong answer against the right one}\par
\vspace{1.5mm}
\widefig{146mm}{fig1_bluff.pdf}\par
\vspace{0.5mm}
\fignote{Registered test passed for \hOnePass{} of \hOneJudges{} judges. Control added after the registered runs, confirmed later.}
\end{frame}

% === 5. The judges know ===
\begin{frame}{The judges can tell that the answer is wrong}
\vspace{1.5mm}
\begin{tabular}{@{}>{\raggedright\arraybackslash}p{33mm}>{\raggedright\arraybackslash}p{102mm}@{}}
\lead{\flaggedMin{} to \flaggedMax{}} & of wrong answers are called incorrect by the judge when asked about one response \\[2mm]
\lead{\bFlaggedOnlyMin{} to \bFlaggedOnlyMax{}} & of duels still go to the wrong answer on those questions, the registered check \\[2mm]
\lead{\bFlaggedMin{} to \bFlaggedMax{}} & where the judge also calls the right answer correct, added after the registered runs \\
\end{tabular}\par
\vspace{4mm}
\begin{exhibit}
\raggedright ``Assistant A provides an answer, but it is incorrect [\ldots] Assistant B does not answer the question. Despite A's factual error, it is more responsive than B, which offers no answer''\par
\vspace{1mm}
\fignote{GPT-5.6 Luna on the duel of slide 2, asked to explain by the MT-Bench prompt. It chose A.}
\end{exhibit}
\end{frame}

% === 6. Registered replications ===
\begin{frame}{It replicates in two more registered question sets}
\widefig{122mm}{fig2_replicate.pdf}\par
\vspace{0.5mm}
{\fontsize{12}{14.5}\selectfont Fresh questions reproduced both rates for \lead{\cfAllPass{} of \cfKept{}} judges. On free-form questions \lead{\ffAboveHalf{} of \ffKept{}} stayed above one half, \ffBelowNames{} the exception at \ffBelowB{}. Rows with $\dagger$ are two larger judges, registered later.\par}
\end{frame}

% === 7. The judge prompt of AlpacaEval 2.0 (Amendment 5, registered before any call with it) ===
\begin{frame}{Same result with the AlpacaEval 2.0 judge prompt}
\vspace{0.5mm}
\stat{\herosmall{\aeBEightMin{}--\aeBEightMax{}}}{40mm}{\keysquare\ wrong answer against ``I don't know'', eight judges}\hfill
\stat{\herosmalldark{\aeCtrlEightMin{}--\aeCtrlEightMax{}}}{38mm}{\keydot\ same wrong answer against the right one}\par
\vspace{2mm}
\begin{exhibit}
\raggedright{\fontsize{12}{14.5}\selectfont ``Your task is to assess these responses, and select the model that produces the best output from a human perspective. [\ldots] make sure your output only contains one of the following model identifiers and nothing else [\ldots]: m or M.''\par}
\vspace{0.5mm}
\fignote{The template of the default AlpacaEval 2.0 judge, used verbatim on the 150 fresh questions.}
\end{exhibit}
\vspace{2mm}
{\fontsize{11.5}{13.5}\selectfont Registered tests A1 and A2 passed for \lead{\aeAOnePass{} of \aeKept{}} judges, A3 for \aeAThreePass{} of \aeLargeKept{} larger ones. We read the written identifier, not its probability, and our judges stand in for GPT-4 Turbo.\par}
\end{frame}

% === 8. What it does to a leaderboard ===
\begin{frame}{Guessing buys Elo points under every judge}
\vspace{1mm}
\widefig{148mm}{fig3_guess.pdf}\par
\vspace{2.5mm}
{\fontsize{12}{14.5}\selectfont Computed from the measured rates. Guessing on the fifth of questions a model does not know gains it \lead{\guessEloMin{} to \guessEloMax{} Elo points} over abstaining.\par}
\vspace{2mm}
\fignote{\sym{W}(0.25) = \sym{b} + 0.25\,(\sym{a} \textminus{} \sym{b}), with \sym{a} and \sym{b} measured on all questions of the registered study.}
\end{frame}

% === 9. What fixes it ===
\begin{frame}{Verify-then-compare brings all six below one half}
\widefig{140mm}{fig4_fixes.pdf}\par
\vspace{1mm}
{\fontsize{12}{14.5}\selectfont \lead{Verify-then-compare} brings all \vtcBelow{} below one half, at a cost in right answers. Instructions manage at most \instrBelowMax{}. Explanations help on free-form questions.\par}
\end{frame}

% === 10. Larger judges and reward models (Amendment 4, registered before any call on these question sets) ===
\begin{frame}{Larger judges and reward models}
\vspace{1mm}
\begin{columns}[T,onlytextwidth]
\begin{column}{0.47\textwidth}
\textbf{Two larger judges}\par
{\fontsize{10.5}{12.5}\selectfont\grey{Registered before any call on these questions}\par}
\vspace{3mm}
\stat{\herosmall{\lgBQwen{}}}{46mm}{\keysquare\ Qwen 3.5 397B, wrong answer against ``I don't know''}\par
\vspace{2mm}
\stat{\herosmall{\lgBGLM{}}}{46mm}{\keysquare\ GLM 5.2, the same}\par
\vspace{2mm}
\stat{\herosmalldark{\lgCtrlMin{}--\lgCtrlMax{}}}{38mm}{\keydot\ same wrong answer against the right one}\par
\vspace{3mm}
{\fontsize{10.5}{12.5}\selectfont On multiple choice L1 and L2 passed for both. On free-form questions the bluff rates were \lgFfBQwen{} and \lgFfBGLM{}, and L3, which needed one judge, \outcome{\lgLThreeOutcome{}} with Qwen.\par}
\end{column}
\hfill
\begin{column}{0.47\textwidth}
\textbf{Three reward models}\par
{\fontsize{10.5}{12.5}\selectfont\grey{Each response scored alone}\par}
\vspace{3mm}
\stat{\herosmall{\rmMcBSepMin{}--\rmMcBSepMax{}}}{36mm}{multiple choice, wrong answer scored above ``I don't know''}\par
\vspace{2mm}
\stat{\herosmalldark{\rmFfBSepMin{}--\rmFfBSepMax{}}}{36mm}{free-form questions, the same}\par
\vspace{3mm}
{\fontsize{10.5}{12.5}\selectfont Rates are over the questions on which a model scores the right answer above the wrong one. R1 on multiple choice passed for \rmROnePass{} of \rmCount{}. R2 on free-form questions \outcome{\rmRTwoOutcome{}}, with \rmRTwoPass{} of \rmCount{} where two were needed.\par}
\end{column}
\end{columns}
\end{frame}

% === 11. Scope and what to do ===
\begin{frame}{Report the bluff rate next to any win rate}
\begin{tabular}{@{}>{\bfseries}p{28mm}>{\raggedright\arraybackslash}p{107mm}@{}}
Scope & A bare verdict on short factual questions, as the default AlpacaEval 2.0 judge reads. \\[1mm]
 & No human raters. \\[0.5mm]
 & A prompt that asks for an explanation and allows a tie shrinks the effect. \\[1mm]
What to do & Measure the bluff rate of the judge and the instruction in use, on about 100 questions with a key, and report it next to any win rate. \\[1mm]
 & Or use verify-then-compare and report how many right answers it costs. \\
\end{tabular}\par
\vspace{1mm}
\url{https://github.com/123yaroslav/win-rate-pays-for-guessing}
\end{frame}

\end{document}
